\section{Introduction}\label{sec:introduction}
This article establishes the rank-four part of the classification program
without a restriction on the fusion multiplicities and without assuming
a pivotal structure. Its conclusion concerns both the based rings that
occur and the sphericality of every category realizing them. The proof
uses one common framework---sphericalization, central induction and exact
character arithmetic---to reduce arbitrary rank-four categories to a
finite, explicitly bounded collection. The exact certificate stage then
identifies the fifteen realizable rings.

The companion finite census article, \emph{Rank-four Grothendieck rings
through multiplicity 128}, gives a larger independent finite verification.
The present proof includes its own smaller census and does not require
the larger census as an input.

\begin{theorem}[Classification]\label{thm:classification}
There are exactly fifteen Grothendieck rings of complex fusion categories
of rank four, up to based isomorphism. They are listed in
Table~\ref{finite:tab:models}. Every complex categorification of each of
these rings admits a spherical structure. Moreover,
\[
 \FPdim(\C)\le24+12\sqrt3,\qquad
 \max_{X,Y,Z}N_{XY}^{Z}\le6.
\]
Both bounds are sharp, and equality in either holds precisely for the
near-group ring $\operatorname{NG}(C_3,6)$.
\end{theorem}

% Generated by supplementary/check_manuscript.py from verified model data.
\begin{table}[htbp]
\centering\small
\begin{tabular}{@{}clcl@{}}
\toprule
$\mu$ & A unitary realization & $D$ & Reference \\
\midrule
$1$ & $\mathrm{Vec}_{C_2\times C_2}$ & $4$ & \cite[\S3.4]{LPR} \\
$1$ & $\mathrm{Fib}\boxtimes\mathrm{Vec}_{C_2}$ & $5+\sqrt{5}$ & \cite[\S3.4]{LPR} \\
$1$ & $\operatorname{Rep}(D_{10})$ & $10$ & \cite[\S3.4]{LPR} \\
$1$ & $\operatorname{Ad}(\mathrm{SU}(2)_6)$ & $8+4\sqrt{2}$ & \cite[Theorem 3.1]{EIP} \\
$1$ & $\mathrm{Fib}\boxtimes\mathrm{Fib}$ & $(15+5\sqrt{5})/2$ & \cite[\S3.4]{LPR} \\
$1$ & $\operatorname{Ad}(\mathrm{SU}(2)_7)$ & $\delta_7$ & \cite[\S3.4]{LPR} \\
$1$ & $\mathrm{Vec}_{C_4}$ & $4$ & \cite[\S3.4]{LPR} \\
$1$ & Tambara--Yamagami on $C_3$ & $6$ & \cite{TambaraYamagami} \\
$1$ & Principal even part of $S'$ & $8+4\sqrt{2}$ & \cite[Theorem 1]{LMP} \\
$2$ & Principal even part of $2D2$ & $20+8\sqrt{5}$ & \cite[Theorem 3.1]{EIP} \\
$2$ & Haagerup category $\mathcal H_1$ & $(39+9\sqrt{13})/2$ & \cite[Table 2]{GrossmanSnyder} \\
$2$ & $\operatorname{Rep}(A_4)$ & $12$ & \cite[Theorem 1.1]{Larson} \\
$3$ & $\operatorname{Ad}(E_8)$ & $\delta_E$ & \cite{IzumiE8,EvansJones} \\
$3$ & Near-group type $(C_3,3)$ & $(21+3\sqrt{21})/2$ & \cite[Example 9.4]{Izumi} \\
$6$ & Near-group type $(C_3,6)$ & $24+12\sqrt{3}$ & \cite[Theorem 10.19]{Izumi} \\
\bottomrule
\end{tabular}
\caption{One unitary realization of each surviving based ring, with its
multiplicity $\mu$ and global Frobenius--Perron dimension $D$.
Here $D_{10}$ has order ten, $S'$ has index $3+2\sqrt2$, and the
near-group category of type $(C_3,3)$ is the Izumi--Xu category.}
\label{finite:tab:models}
\smallskip
\begin{minipage}{\textwidth}\footnotesize
The numbers $\delta_7$ and $\delta_E$ are the largest real roots of
$x^3-27x^2+162x-243$ and
$x^4-45x^3+375x^2-1125x+1125$, respectively.
The supplementary model check verifies these polynomials, the positive
dimension characters, and all tensor identifications exactly.
\end{minipage}
\end{table}


The classification and sphericality are proved in
Corollary~\ref{finite:cor:classification}, and sharpness in
Corollary~\ref{finite:cor:sharp-bounds}.
The classification is computer assisted. The complete local census below
$D<1000$ contains $\CensusTotal$ fusion rings: $\CensusSelfDual$ self-dual
rings and $\CensusNonSelfDual$ rings with a dual pair.
The supplementary package supplies its own enumeration and classification
programs, every full multiplication tensor, and the exact certificate
ledger. Table~\ref{finite:tab:filters} records the successive reductions
to the fifteen models. The finite verification is
specified in Section~\ref{finite:sec:census}, and all its mathematical
exclusion criteria are established in Section~\ref{sec:finite-criteria}.
In particular, positive-dimensional indicator estimates are invoked only
after a proved pseudo-unitary reduction. Other branches are tested through
the sphericalization, whose signed quotient and dimension character are
both exhaustively retained.

The bound that makes this finite verification exhaustive is independent
of the computation.
\begin{theorem}[Uniform bound]\label{thm:uniform-global}
For every complex fusion category $\C$ of rank four,
\[
 \FPdim(\C)<1000,\qquad
 N_{XY}^{Z}\le31
 \quad(X,Y,Z\in\operatorname{Irr}\C).
\]
\end{theorem}

The bound is independent of the fusion-ring census. Its constants come
from the dimension and character estimates proved below; no sharpness
is asserted for this theoretical reduction.
Its proof is completed in Section~\ref{sec:global-bound-proof}.
Central induction in the sphericalization gives positive integral
Galois columns and a descent theorem for the dimension field.
Quadratic arithmetic and averaged indicator identities then give
$D<1000$ in the nonspherical quadratic case and $D<950$ in the
nonspherical quartic case. Pivotal categories satisfy the independent
bound $D<1000$.

The structural ingredients also prove sphericality when the category
has a non-self-dual simple, a nontrivial invertible simple, or a cubic Frobenius--Perron
character field. The two Fibonacci normal forms needed in the proof
are treated directly. A short
quadratic argument treats the Haagerup ring. These results establish
sphericality of every categorification of each surviving ring; existence
of one spherical realization alone would not suffice for that conclusion.

A character field and its integral based order are distinguished
throughout. In particular, a quadratic character component may have a
constant rational formal codegree, and this possibility is retained in
the signed reductions. The arithmetic charts used for the uniform bound
are bounded integer calculations derived within the proof, rather than
assumptions taken from a fusion-ring census.

\subsection*{Notation and proof structure}
The ordinary and signed dimension characters are denoted by $d_i,x_i$,
with global dimensions $D=\sum d_i^2$ and $E=\sum x_i^2$.
The ordinary and signed formal codegrees are $f_i,g_i$;
$u=\sum f_i^{-2}$ and $v=\sum g_i^{-2}$.
A character-value field is not silently identified with the field of
its global codegree: that identification is supplied where needed by
Section~\ref{sec:central-columns}. The proof separates as follows.
\begin{center}\small
\begin{tabular}{p{.31\textwidth}p{.59\textwidth}}
\toprule
Case&Reduction used in the main proof\\
\midrule
Non-self-dual rings&Automatic sphericality; a theoretical candidate
classification and exclusion (Sections~\ref{sec:ns}, \ref{finite:sec:catalog}).\\
Self-dual, quadratic codegree&Dimension-field descent and quadratic
arithmetic (Section~\ref{sec:quad}).\\
Cubic character field&Positive central measure and periodic moments
(Section~\ref{sec:cub}).\\
Quartic global codegree&Central ratios and a bounded arithmetic chart
(Section~\ref{sec:quartic-bound}).\\
Spherical case&Independent pivotal bound (Section~\ref{sec:pivotal-bound}).\\
Bounded tables&Complete ordinary census, signed-cover verification,
and exact certificates (Section~\ref{finite:sec:catalog}).\\
\bottomrule
\end{tabular}
\end{center}
All categorical reductions needed for the classification are proved in
this article. The exact enumeration and certificates are supplied in the
\texttt{unrestricted} package, with their interfaces described in the
verification guide. Existence of the named unitary categories is supplied
by the construction references in Table~\ref{finite:tab:models}.

